\section{The complete small spectrum and a three-scale family}\label{sec:spectrum}

The linearization of HYM at the parallel metric is the heat
equation for the induced connection Laplacian on trace-free
Hermitian endomorphisms. Write $\mathfrak p_r=\Herm_0(r)$,
$d=r^2-1$, $m=4d$, and
\[
 \mathcal L_{\tau,\alpha,U}=D^*D,\qquad
 \mathcal E_\tau(s)=\int_{C_\tau}|Ds|^2\,dA_\tau .
\]
In this section the unitary transports may vary between fibres.
For $X\in\mathfrak p_r^4$ set
\begin{equation}\label{eq:effective-forms}
 \begin{split}
 M(X)&=\sum_i A_i\|X_i\|^2,\\
 g_\alpha(X)&=\sum_{i=+,-}\alpha_i^2\kappa_i\|[T_i,X_i]\|^2,\\
 Q(X)&=g_\alpha(X)+
       \sum_{e=ij}w_e\|X_i-\Ad(U_{ij})X_j\|^2,\qquad
 w_e=2\pi/L_e .
 \end{split}
\end{equation}
These are respectively the mass form and the Hessian of
\eqref{eq:potential} at $o$. They are defined by the component
geometry, harmonic forms and node transports before solving a
global analytic problem.

\begin{theorem}\label{thm:relative-spectrum}
Let $\nu_0\le\cdots\le\nu_{m-1}$ be the generalized eigenvalues
of $(Q,M)$, including zeros, and let
$\lambda_0\le\lambda_1\le\cdots$ be those of
$\mathcal L_{\tau,\alpha,U}$. There are $\gamma>0$ and
$e(\alpha_*,b)\to0$ as $\alpha_*,b\to0$ such that
\begin{equation}\label{eq:relative-spectrum}
 (1-e)\nu_k\le\lambda_k\le(1+e)\nu_k\quad(0\le k<m),
 \qquad \lambda_m\ge\gamma.
\end{equation}
Exactly $m$ eigenvalues, counting zeros, tend to zero. The kernels
of the finite and actual operators correspond exactly through
\eqref{eq:common-kernel}. The estimates are uniform in all
$\alpha_i^2L_e$, all $L_e/L_f$, and all unitary transports,
including reducible ones.

Component averaging identifies the $m$-dimensional spectral
subspace associated with the eigenvalues tending to zero
asymptotically isometrically with $(\mathfrak p_r^4,M)$.
Let $\Pi$ be the orthogonal projection onto this small-eigenvalue
subspace, and let $P_a$ be the orthogonal projection onto the
image of the trial map $J_a$ constructed below. Then
\begin{equation}\label{eq:projection-comparison}
 \|\Pi-P_a\|_{L^2\to L^2}
                \le C\sqrt{\alpha_*^2+w_*},\qquad
 w_*=\max_e w_e.
\end{equation}
On the complementary space the inverse, heat and resolvent
bounds are respectively $\gamma^{-1}$, $e^{-\gamma s}$
and $(\gamma+\operatorname{Re}z)^{-1}$ for
$\operatorname{Re}z\ge0$.
\end{theorem}

The scalar graph approximation on degenerating surfaces has a
substantial history, including relative estimates with arbitrary
short-edge ratios \cite{burger}. Here the point requiring proof
is the combined component-and-node estimate, including the
commutator kernel and every weaker positive scale.

\subsection{Trial sections}
Choose a small fixed radius $a$ in the branch disks, independent
of all degenerating parameters. On an elliptic component begin
with the constant matrix $X_i$ in the harmonic gauge. On
$a<|u|<2a$ interpolate it with its parallel extension having
node value $X_i$. The primitive of the harmonic one-form obeys
$|F_{ie}|\le Ca$ and $|dF_{ie}|\le C$. The exact bound
\[
 \|\Ad(e^{i\alpha_iT_iF_{ie}})X_i-X_i\|
           \le|\alpha_iF_{ie}|\|[T_i,X_i]\|
\]
retains the component energy factor in every cutoff term.
The component energy is at most $(1+Ca^2)g_\alpha(X)$.
On rational components use constants. Conformal invariance
transfers these energies to the fibre exactly.

On the remaining node annulus, whose conformal length is
$\ell_e=L_e+2\log a$, interpolate the two parallel-frame
values affinely in logarithmic radius. Its energy is exactly
$(2\pi/\ell_e)\|X_i-\Ad(U_e)X_j\|^2$. Denote the resulting
continuous $H^1$ trial section by $J_aX$. Then
\begin{equation}\label{eq:trial-estimates}
 \begin{split}
 \mathcal E_\tau(J_aX)&\le(1+Ca^2)g_\alpha(X)
       +\sum_e\frac{2\pi}{\ell_e}
                      \|X_i-\Ad(U_e)X_j\|^2,\\
 1\le L_e/\ell_e&\le1+C_a/L_*,\\
 \bigl|\|J_aX\|_{L^2}^2-M(X)\bigr|
       &\le(Ca^2+o_a(b))M(X).
 \end{split}
\end{equation}
The mass estimate uses the $O(a^2)$ disk and annular areas,
unitarity and bounded interpolation, and smooth metric
convergence on fixed component exteriors. It is uniform in
the transports and length ratios. Min--max gives the relative
upper bound in \eqref{eq:relative-spectrum}, first choosing
$a$ small and then the parameters small. For one fixed $a$
it also gives $\lambda_{m-1}\le C(\alpha_*^2+w_*)$.

\subsection{Component means with an energy-controlled frame error}
On the punctured elliptic domain $\Omega_i(a)$, let
$\vartheta_{i,a}$ be the absolute harmonic representative of
the restriction class of $\vartheta_i$, with zero normal
component at the boundary. Write
$\vartheta_i=\vartheta_{i,a}+df_{i,a}$ and
$\kappa_{i,a}=\|\vartheta_{i,a}\|_{L^2}^2$.
Then $\kappa_{i,a}\to\kappa_i>0$. For completeness, the
restriction of $\vartheta_i$ is a competitor for the energy
minimization, giving a uniform upper bound. A weakly
convergent harmonic subsequence on compact subsets has finite
$L^2$ energy across every puncture. Its holomorphic
representation has no negative Laurent terms with finite
$L^2$ norm, so it extends across the puncture. Its periods,
including zero puncture periods, are those of
$\vartheta_i$. Uniqueness of the harmonic representative
on the closed component and lower semicontinuity prove
the asserted norm convergence.

Gauge by $e^{i\alpha_iT_if_{i,a}}$ and diagonalize $T_i$.
Each matrix entry then has connection
$d-i\beta\vartheta_{i,a}$ with
$\beta=\alpha_i(t_b-t_c)$. Decompose an entry as $v=z+w$,
where $w$ has mean zero. Coclosedness and the zero normal
condition give
$\langle dw,\vartheta_{i,a}z\rangle=0$.
The fixed-domain Poincar\'e inequality and a small-$\beta$
perturbation give
$\|w\|\le C_a\|(d-i\beta\vartheta_{i,a})w\|$.
The remaining cross term is bounded by
$C_a\beta^2|z|\|w\|$. Completing the square therefore yields
\begin{equation}\label{eq:component-lower-entry}
 \|(d-i\beta\vartheta_{i,a})v\|^2
 \ge\frac12\|(d-i\beta\vartheta_{i,a})w\|^2
            +(\beta^2\kappa_{i,a}-C_a\beta^4)|z|^2.
\end{equation}
Summing matrix entries defines a real linear average
$s\mapsto X_i(s)$ and gives
\begin{equation}\label{eq:component-lower}
 \begin{split}
 \mathcal E_i(s)&\ge
   (1-C_a\alpha_*^2)\alpha_i^2\kappa_{i,a}
                                  \|[T_i,X_i]\|^2,\\
 \|w_i\|_{H^1}^2+\alpha_i^2\|[T_i,X_i]\|^2
                                      &\le C_a\mathcal E_i(s).
 \end{split}
\end{equation}
No division by $\alpha_i$ or an eigenvalue difference is used.
On rational components take the ordinary mean. These estimates
concern actual component energy by conformal invariance.

Let $c_{ie}$ be the boundary-circle mean in the parallel node
frame. The change from the averaging gauge to this frame
acting on $X_i$ is at most
$C_a|\alpha_i|\|[T_i,X_i]\|$. The remaining error is bounded
by the $H^1$ trace norm of $w_i$. Thus
\begin{equation}\label{eq:boundary-means}
 \|c_{ie}-X_i(s)\|\le C_a\mathcal E_i(s)^{1/2}.
\end{equation}
Replacing this by an error proportional to
$|\alpha_i|\|X_i\|$ would lose relative control near a
changing commutant.

\subsection{The combined lower form}
Partition the energy into the disjoint truncated components
and their complementary node annuli:
$\mathcal E_\tau=\mathcal E_{\mathrm{comp}}
                  +\mathcal E_{\mathrm{neck}}$.
Angular averaging and the one-dimensional Dirichlet principle
give
\[
 \sum_e\frac{2\pi}{\ell_e}
       \|c_{ie}-\Ad(U_e)c_{je}\|^2
                   \le\mathcal E_{\mathrm{neck}}.
\]
Apply the triangle inequality in the direct sum of edge
spaces weighted by $\sqrt{w_e}$. Using
$w_e\le2\pi/\ell_e$ and \eqref{eq:boundary-means} gives
\begin{equation}\label{eq:weighted-neck-lower}
 \left(\sum_e w_e\|X_i-\Ad(U_e)X_j\|^2\right)^{1/2}
 \le\mathcal E_{\mathrm{neck}}^{1/2}
           +C_a\sqrt{w_*}\,\mathcal E_{\mathrm{comp}}^{1/2}.
\end{equation}
With $\delta=C_a\sqrt{w_*}$ use
$(x+\delta y)^2\le(1+\delta)x^2+(\delta+\delta^2)y^2$.
Meanwhile \eqref{eq:component-lower} and
$\kappa_{i,a}\to\kappa_i$ give
$g_\alpha(X)\le(1+\theta(a)+C_a\alpha_*^2)
                   \mathcal E_{\mathrm{comp}}$, where
$\theta(a)\to0$. Adding the two estimates yields
\begin{equation}\label{eq:combined-form-lower}
 Q(X(s))\le
   \{1+\theta(a)+C_a\alpha_*^2+O_a(\sqrt{w_*})\}
                                    \mathcal E_\tau(s).
\end{equation}
The component and annular energies have not been counted
twice. This estimate is relative to their actual sum and
contains no smallest scale in a denominator.

\subsection{Mass and completeness}
We give the mass estimate needed both for the lower eigenvalue
bound and for completeness of the mode list. Let $s$ be a unit
eigensection with eigenvalue $\lambda\le1$. Equation
\eqref{eq:component-lower} bounds its mean-zero component
parts in $H^1$ by $C_a\sqrt\lambda$. On fixed collars,
interior elliptic estimates improve the boundary error to
the same order in supremum norm. In the equation for the
mean-zero part, the Laplacian of the gauged constant is
$\alpha_i^2|\vartheta_{i,a}|^2\operatorname{ad}(T_i)^2X_i$.
Its norm is bounded by $C_a|\alpha_i|\sqrt{\mathcal E_i}$
using \eqref{eq:component-lower}. The other inhomogeneous
term is $\lambda s$. Enlarge the fixed collar across the
circle before applying the interior estimate. In the
actual conformal metric the connection Laplacian is
multiplied by a bounded smooth scalar there, without a new
first-order term. The leading boundary values have size
at most $C\|X_i\|$ with $C$ independent of $a$, because
the changes of frame are unitary.

On each node annulus split the section into its entrywise
harmonic Dirichlet extension $h$ and a zero-boundary remainder
$v$. The maximum principle and annular area $O(a^2)$ bound
the mass of $h$ by $Ca^2+C_a\sqrt\lambda$. Orthogonality
of Dirichlet energies gives
$\mathcal E(v)\le\mathcal E(s)$. The zero-boundary Poincar\'e
estimate
\begin{equation}\label{eq:annulus-poincare}
 \|v\|_{L^2(dA_\tau)}^2\le Ca^2\mathcal E(v)
\end{equation}
is uniform in the annular modulus. To verify it, the
metric is bounded by a fixed multiple of
$|du|^2+|dv_{\mathrm{branch}}|^2$, with
$v_{\mathrm{branch}}=\tau_e/u$. For the first branch,
the zero outer boundary value and radial Cauchy--Schwarz
give
$|v(r,\theta)|^2\le\log(a/r)\int_r^a
|\partial_tv(t,\theta)|^2t\,dt$.
Integration against $r\,dr\,d\theta$ is bounded by
$Ca^2$ times the radial energy. The other branch gives
the same estimate for its area term, using the zero
boundary value at its outer circle. This proves
\eqref{eq:annulus-poincare}.

The component mean decomposition, metric convergence on
fixed exteriors, and these annular bounds now imply that
on any span of at most $m+1$ orthonormal eigensections
with eigenvalues at most $\Lambda\le1$,
\begin{equation}\label{eq:weak-mass}
 |M(X(s))-\|s\|_{L^2}^2|
 \le\{Ca^2+C_a\sqrt\Lambda+o_a(b)\}\|s\|_{L^2}^2.
\end{equation}
Passing from individual eigensections to this span costs
only its fixed dimension. The leading omitted-area
constant is independent of $a$.

If $\lambda_m$ tended to zero along a joint degeneration,
first choose $a$ small and then take the parameters far
along that sequence. Equation \eqref{eq:weak-mass} would
make averaging injective on an $(m+1)$-dimensional space
into $\mathfrak p_r^4$, a contradiction. This proves
the uniform complementary gap. The trial upper bound
already puts $m$ eigenvalues below it.

For the first $k+1$ eigensections use
$\Lambda=C(\alpha_*^2+w_*)$ in \eqref{eq:weak-mass}.
Their averages form a $(k+1)$-dimensional space.
Finite min--max gives a section in this span with
$Q(X(s))\ge\nu_kM(X(s))$. Combining this with
\eqref{eq:combined-form-lower} proves the relative lower
bound. For any requested error first choose $a$, then
choose $\alpha_*,b$ independently small. These quantifiers
give the function $e(\alpha_*,b)\to0$ in the theorem,
uniformly in all remaining ratios and transports.

The zero set of $Q$ is exactly
\eqref{eq:common-kernel}, since every node weight is
positive. Section~\ref{sec:geometry} identifies these
tuples with all global parallel endomorphisms, which are
precisely the kernel of $D^*D$. Thus zeros agree exactly
through all kernel changes.

Choose $a=a(\alpha_*,b)$ tending sufficiently slowly to
zero that the preceding errors vanish and
$|\log a|/L_*\to0$. The averaging and trial mass
estimates are then asymptotic isometries. On the
complement of the small-eigenvalue subspace, the gap gives
\[
 \|(1-\Pi)J_aX\|^2\le
       \gamma^{-1}\mathcal E_\tau(J_aX)
       \le C(\alpha_*^2+w_*)M(X).
\]
The spaces have equal dimension and the trial mass
has a uniform lower bound. Equality of their maximal
principal angle with the norm of the difference of
orthogonal projections proves
\eqref{eq:projection-comparison}. The spectral theorem
gives the stated inverse, heat and resolvent estimates
on the complement. On the positive small-eigenvalue subspace the
smallest rate and inverse norm instead have the
relative comparison with the smallest positive
$\nu_k$; they need not stay bounded away from zero
and infinity. This completes the proof of
Theorem~\ref{thm:relative-spectrum}.

\subsection{A three-scale family}
Order the vertices as $E_+,R_p,E_-,R_q$, numbered $0,1,2,3$.
Use identity transports on $0$--$1$--$2$--$3$ and
$U=(I+i\sigma_x)/\sqrt2$ on $3$--$0$. Set $T_+=\sigma_z$,
$\alpha_-=0$, and, for $h>0$, take
\begin{equation}\label{eq:h-family}
 w_{01}=w_{12}=w_{23}=h,\quad w_{30}=h^2,\quad
 \alpha_+=h^2/\sqrt{\kappa_+},\quad
 \tau_e=\exp(-2\pi/w_e).
\end{equation}
These are genuine parameters of the compact family.
Section~\ref{sec:geometry} proves stability for all
sufficiently small $h>0$.

In the Hilbert--Schmidt orthonormal Pauli basis
$(\sigma_x,\sigma_y,\sigma_z)/\sqrt2$ at each vertex,
$M=I$ and the matrix of $Q$ is
\begin{equation}\label{eq:h-matrix}
 Q_h=h(L_{P_4}\otimes I_3)+h^2D_{30}^TD_{30}
          +h^4\operatorname{diag}(4,4,0,0,\ldots,0),
\end{equation}
where
\[
 L_{P_4}=\begin{pmatrix}
 1&-1&0&0\\-1&2&-1&0\\0&-1&2&-1\\0&0&-1&1
 \end{pmatrix},\qquad
 D_{30}X=X_3-RX_0,\qquad
 R=\begin{pmatrix}1&0&0\\0&0&1\\0&-1&0\end{pmatrix}.
\]
The positive path eigenvalues are $2-\sqrt2,2,2+\sqrt2$.
Thus nine finite eigenvalues are these three coefficients
times $h$, each with multiplicity three, up to $O(h^2)$.
The isometric constant-tuple inclusion
$Z\mapsto(Z,Z,Z,Z)/2$ compresses the cycle term to
$\operatorname{diag}(0,\tfrac12,\tfrac12)$ and the
component term to $\operatorname{diag}(1,1,0)$.
Schur complementation across the order-$h$ complement
therefore gives two eigenvalues
$\tfrac12h^2+O(h^3)$.

For the last eigenvalue one must use its invariant sector,
rather than this absolute remainder. The $\sigma_x$
sector is invariant, its graph operator has exact
constant kernel, and its positive gap is at least
$(2-\sqrt2)h$. The order-$h^4$ perturbation has
expectation $h^4$ on its unit constant vector.
The complementary inverse bound gives the last
eigenvalue $h^4+O(h^7)$. As a check on the displayed
finite matrix, its determinant is
\begin{equation}\label{eq:h-determinant}
 \det Q_h=16h^{17}(1+3h)(1+2h^2+6h^3).
\end{equation}
All entries and the compressions used in this calculation
are printed above.

\begin{corollary}\label{cor:three-scales}
On the stable compact family \eqref{eq:h-family}, nine
actual eigenvalues have leading coefficients
$2-\sqrt2,2,2+\sqrt2$, each three times, at order $h$;
two are asymptotic to $\tfrac12h^2$; one is
asymptotic to $h^4$. The next eigenvalue has a uniform
positive lower bound. On the fixed small basin of
Theorem~\ref{thm:all-time}, actual nonlinear HYM
trajectories are approximated by
\eqref{eq:finite-flow} through time $h^{-4}$ and all
later times, with absolute integrated squared-distance
error at most
$\eta+O(h(1+|\log h|)^2)$. The spatial initial class
of Theorem~\ref{thm:spatial-entry} has the corresponding
all-time conclusion after preparation.
\end{corollary}

\begin{proof}
Apply Theorem~\ref{thm:relative-spectrum} to the finite
eigenvalue estimates. Apply Theorems~\ref{thm:all-time}
and \ref{thm:spatial-entry} separately to the same
fixed matrices and transports. Here
$\epsilon_*=h/(2\pi)$, $\alpha_*=h^2/\sqrt{\kappa_+}$,
and $b=e^{-2\pi/h}$, so
$a_*=O(h(1+|\log h|)^2)$. The nonlinear conclusion
follows from the residual and energy proof, not from
the eigenvalue asymptotics.
\end{proof}
